% Overleaf: Compiler = LuaLaTeX
% 日本語: Menu → Settings → Compiler → LuaLaTeX
\documentclass[
  a4paper,
  11pt,
  dvipdfmx,
  titlepage
]{jlreq}

\usepackage[margin=25mm]{geometry}
\usepackage{graphicx}
\usepackage{amsmath,amssymb}
\usepackage{booktabs}
\usepackage{multirow}
\usepackage{hyperref}
\usepackage{cleveref}
\usepackage{enumitem}
\usepackage{siunitx}
\sisetup{detect-all=true}

\hypersetup{
  colorlinks=true,
  linkcolor=blue,
  citecolor=blue,
  urlcolor=blue
}

\title{%
  工場内AGV環境における危険度連動経路可視化eHMIの\\%
  設計とVR実験計画
}
\author{%
  著者名 \\%
  \small 所属・メールアドレス
}
\date{\today}

\begin{document}
\maketitle

\begin{abstract}
  スマートファクトリーにおいて，AGV（Automated Guided Vehicle）と作業員の共用空間は拡大しているが，
  多数のAGVが同時走行する環境では，作業員が衝突リスクを予測しにくいという課題がある．
  本研究は，AR風の経路オーバーレイを用いたeHMI（external Human-Machine Interface）により，
  AGVの残り走行経路を提示し，\textbf{危険度に応じて色と不透明度を動的に変化}させる提案手法（Proposed）を設計した．
  危険度は，参加者の視線方向に基づく動的基準線とAGV残存経路の交差から得られる
  Time-to-Conflict（TTC），および近接距離の2指標を統合して算出する．
  提案手法の有効性を検証するため，Meta Quest 3を用いたVR実験を計画する．
  被験者内3条件（Baseline / No-AR / Proposed）×10ケースの構成とし，
  完了時間，ニアミス回数，移動距離を従属変数とする．
  同一ケース番号ではAGVの台数・速度・軌道をシード固定することで条件間の比較妥当性を確保する．
\end{abstract}

\tableofcontents
\newpage

\section{はじめに}
\label{sec:intro}

\subsection{背景}

製造現場の自動化が進む中，AGVは荷物搬送の中核として導入されている．
AGV導入により搬送効率は向上する一方，作業員とAGVが同一通路を共有する場面が増え，
\textbf{衝突・ニアミス・歩行の中断}といった安全上のリスクが顕在化している．

従来，AGVの意図は音，光，表示板などのeHMIで伝達されるが，
\textbf{多数台が同時に存在する}工場環境では，どのAGVが自分にとって危険かを
瞬時に判断することが困難である．
特に作業員が目的地へ移動しながら周囲を監視するタスクでは，
全AGVの経路を等しく表示するだけでは情報過多となり，
注意資源の配分が非効率になる可能性がある．

\subsection{課題設定}

本研究が扱う中心的課題は次の通りである．

\begin{quote}
  \textbf{研究課題（RQ）:}
  AGVの残存走行経路をAR風に可視化する際，
  \emph{危険度に応じた視覚エンコーディング}は，
  作業員の移動効率と安全性（ニアミス）にどのような影響を与えるか？
\end{quote}

この課題に答えるため，以下の3条件を被験者内で比較する．

\begin{itemize}[leftmargin=2em]
  \item \textbf{Baseline}: 全AGVの経路を同一の固定色・不透明度で表示
  \item \textbf{No-AR}: 経路表示なし（情報なしの対照）
  \item \textbf{Proposed}: TTCと近接距離に基づく危険度連動表示（本提案）
\end{itemize}

\subsection{本稿の構成}

\Cref{sec:story}で研究ストーリーと仮説を述べ，
\Cref{sec:method}で提案eHMIの数理モデルを説明する．
\Cref{sec:experiment}でVR実験計画，
\Cref{sec:validity}で実験の妥当性を論じ，
\Cref{sec:implementation}で実装概要を示す．

\section{研究ストーリーと仮説}
\label{sec:story}

\subsection{研究の流れ（ストーリーライン）}

本研究は，次の論理構造に基づいている．

\begin{enumerate}[leftmargin=2em]
  \item \textbf{問題}: 工場内でAGVと人が共存する際，衝突リスクの予測が困難
  \item \textbf{既存アプローチの限界}: 経路を常時・均一に表示するeHMIは，
        危険の程度を区別できず，注意の優先順位付けを支援しない
  \item \textbf{提案}: 残存経路を提示しつつ，
        \emph{相対的危険度}に応じて色相（連続）と不透明度（段階）を変化させる
  \item \textbf{評価}: VR上の工場タスク（Station A$\rightarrow$B移動）で，
        効率（時間・移動距離）と安全（ニアミス）を定量的に比較
  \item \textbf{期待}: Proposedは，No-ARより安全に，Baselineより情報過多を抑えつつ
        適切な注意喚起を実現する
\end{enumerate}

\subsection{理論的根拠}

危険度の時間的予測として\textbf{TTC（Time-to-Conflict）}を用いる点は，
交通・ロボット分野で広く用いられる指標に依拠する\cite{markkula2020}．
また，eHMI研究では，外部表示が歩行者の意図理解や信頼に影響することが報告されている\cite{francis2020eHMI,matthiesen2021survey}．

本提案の新規性は，次の3点に集約される．

\begin{itemize}[leftmargin=2em]
  \item \textbf{視線追従型の動的基準線}:
        目的地最短経路ではなく，参加者が\emph{現在注視している方向}に
        AGVが割り込みうる領域を定義する
  \item \textbf{二重危険度指標の統合}:
        交差予測（TTC）と近接距離の最大値により，
        「交差しないが近い」「交差するがまだ遠い」両ケースをカバー
  \item \textbf{連続色＋段階的不透明度}:
        危険の\emph{連続的変化}（色相）と\emph{警告レベル}（不透明度4段階）を
        分離してエンコード
\end{itemize}

\subsection{研究仮説}

被験者内計画に基づき，以下の仮説を設定する．

\begin{description}[leftmargin=0em,style=nextline]
  \item[\textbf{H1（安全性）}]
    Proposed条件は，No-AR条件と比較して\textbf{ニアミス回数が少ない}．
    根拠: 危険AGVの経路が視覚的に強調され，回避行動が促進される．

  \item[\textbf{H2（効率）}]
    Proposed条件は，Baseline条件と比較して\textbf{完了時間が短い}，
    または\textbf{移動距離が短い}．
    根拠: 危険度に応じた選択的強調により，不要な情報処理負荷が低減される．

  \item[\textbf{H3（情報なし対照）}]
    No-AR条件は，Proposed条件と比較して\textbf{完了時間が長い}，
    または\textbf{移動距離が長い}．
    根拠: 経路情報がないと，AGVの動き予測が困難になり迂回・停止が増える．
\end{description}

\noindent
\textbf{注記:}
H1--H3は独立ではなく，ProposedがBaselineとNo-ARの\emph{中間的な最適点}を
示すことを期待する．
統計検定は，被験者内反復測定のANOVA（または Friedman 検定）と
事後比較（Bonferroni 補正）を想定する．

\section{提案手法：危険度連動eHMI}
\label{sec:method}

\subsection{システム概要}

各AGVに対し，毎フレーム以下の処理を行う（\Cref{fig:system}参照）．

\begin{enumerate}[leftmargin=2em]
  \item 参加者の視線方向から\textbf{動的基準線}（太線ゾーン）を更新
  \item AGVの残存計画経路と基準線の交差から\textbf{TTC}を算出
  \item TTCと近接距離から危険度スコア $R \in [0,1]$ を計算
  \item $R$をHSV色空間の色相・彩度・明度および4段階不透明度に変換
  \item 走行中は残存経路，停止中は停止線として描画
\end{enumerate}

\begin{figure}[htbp]
  \centering
  \fbox{\parbox{0.92\linewidth}{%
    \small
    \textbf{[図1: システム構成（ここに図を挿入）]}\\[0.5em]
    参加者 $\rightarrow$ DynamicCrossingLineTracker（視線基準線）\\
    AGV $\rightarrow$ VehicleRiskCalculator（$R$）\\
    $\rightarrow$ RiskToVisualMapper $\rightarrow$ PathRenderer
  }}
  \caption{提案eHMIの処理パイプライン}
  \label{fig:system}
\end{figure}

\subsection{動的基準線（視線方向の太線ゾーン）}

従来のNavMesh最短経路に基づく基準線ではなく，
参加者が\textbf{現在見ている方向}に注目領域を定義する．

HMDのforwardをXZ平面に投影した単位ベクトル $\mathbf{f}$ に対し，
視線起点 $\mathbf{o}$ からレイキャストを行い，障害物（棚・壁）との
交点 $\mathbf{h}$ を求める．
床面上の基準軸は
\begin{equation}
  \mathbf{a}_\text{start} = \mathbf{p}, \quad
  \mathbf{a}_\text{end} = \mathrm{flatten}(\mathbf{h})
\end{equation}
とする．$\mathbf{p}$ は参加者位置，半幅 $w=\SI{0.5}{m}$（全幅\SI{1.0}{m}）の
帯状ゾーンとして，AGV経路との交差判定に用いる．

この設計により，作業員が立ち止まり別方向を見た場合も，
\textbf{視線に追従した危険評価}が可能となる．

\subsection{危険度スコア}

\subsubsection{表示可否}

まず描画対象か判定する．
\begin{equation}
  \text{visible} = (T \le T_\text{max}) \vee (d \le D_\text{max})
\end{equation}
ここで $T$ はTTC，$d$ はAGVと参加者の3次元距離，
$T_\text{max}=\SI{4.0}{s}$，$D_\text{max}=\SI{13.6}{m}$ である．
表示距離上限は
\begin{equation}
  D_\text{max} = (v_\text{AGV,max} + v_\text{ped}) \times T_\text{max}
  = (2.0 + 1.4) \times 4 = \SI{13.6}{m}
\end{equation}
として，AGV最大速度 $v_\text{AGV,max}=\SI{2.0}{m/s}$，
歩行速度 $v_\text{ped}=\SI{1.4}{m/s}$ から導出する．

\subsubsection{TTC}

AGVの残存計画経路 $\{\mathbf{w}_0,\ldots,\mathbf{w}_{m-1}\}$ 上を
現在速度 $v$ で走行したとき，基準線ゾーンとの最初の交差までの
累積距離 $s$ から
\begin{equation}
  T = \frac{s}{\max(v,\, v_\text{min})}, \quad v_\text{min}=\SI{0.1}{m/s}
\end{equation}
と定義する．交差が $T_\text{max}$ 以内に見つからない場合 $T=\infty$ とする．

\subsubsection{TTC・近接由来スコアと統合}

\begin{align}
  R_\text{ttc} &=
  \begin{cases}
    0 & T = \infty \\
    \left(1 - \dfrac{T}{T_\text{max}}\right)^{\gamma} & 0 \le T \le T_\text{max}
  \end{cases} \\[6pt]
  R_\text{prox} &=
  \begin{cases}
    0 & d \ge D_\text{max} \\
    \left(1 - \dfrac{d}{D_\text{max}}\right)^{\gamma} & d < D_\text{max}
  \end{cases} \\[6pt]
  R &= \max(R_\text{ttc},\, R_\text{prox},\, R_\text{floor})
\end{align}
非線形係数 $\gamma=0.6$，下限 $R_\text{floor}=0.08$ とする．

\subsection{視覚マッピング}

危険度 $R$ をHSV色空間へ写像する（Proposed条件のみ）．
\begin{align}
  H(R) &= (1-R) \times 180° \quad \text{（青緑$\rightarrow$赤）} \\
  S(R) &= 0.4 + 0.6R \\
  V(R) &= 0.5 + 0.3R
\end{align}
不透明度は認知しやすさのため4段階に量子化する．
\begin{equation}
  \alpha(R) \in \{0.35,\, 0.55,\, 0.80,\, 1.00\}
\end{equation}

\subsection{3条件の対照}

\Cref{tab:conditions}に実験条件を示す．
Baselineは「情報あり・危険度なし」，No-ARは「情報なし」，
Proposedは「情報あり・危険度連動」として，被験者内の3水準を構成する．

\begin{table}[htbp]
  \centering
  \caption{実験条件の対照}
  \label{tab:conditions}
  \begin{tabular}{lccc}
    \toprule
    項目 & Baseline & No-AR & Proposed \\
    \midrule
    経路表示 & あり & なし & あり \\
    危険度計算 & なし（$R=1$固定） & --- & TTC + 近接 \\
    色 & 固定青 & --- & HSV連動 \\
    不透明度 & 1.0固定 & --- & 4段階 \\
    表示距離 & \SI{13.6}{m}以内 & --- & \SI{13.6}{m}以内 \\
    \bottomrule
  \end{tabular}
\end{table}

\subsection{経路描画}

走行中は残存計画経路をテーパー付きライン（根元\SI{0.65}{m}，先端\SI{0.12}{m}）で描画する．
停止中（荷物ピックアップ／ドロップ）は停止線（全幅\SI{1.0}{m}）を表示する．
危険度の高いAGVほど手前に描画されるよう，$\mathrm{sortingOrder}=\mathrm{round}(R \times 100)$ とする．

\begin{table}[htbp]
  \centering
  \caption{主要パラメータ一覧}
  \label{tab:params}
  \begin{tabular}{lll}
    \toprule
    記号 & 値 & 意味 \\
    \midrule
    $T_\text{max}$ & \SI{4.0}{s} & TTC表示上限 \\
    $D_\text{max}$ & \SI{13.6}{m} & 近接表示上限 \\
    $\gamma$ & 0.6 & 危険度カーブの急峻さ \\
    $R_\text{floor}$ & 0.08 & 表示中の最低危険度 \\
    $w$ & \SI{0.5}{m} & 基準線半幅 \\
    \bottomrule
  \end{tabular}
\end{table}

\section{VR実験計画}
\label{sec:experiment}

\subsection{実験環境}

Unity（URP）上に構築した工場シーン \texttt{FactoryExperiment} を用い，
Meta Quest 3（OpenXR）で被験者に没入させる．
10〜20台のAGVが荷物のピックアップとドロップを反復する環境で，
被験者はStation AからStation Bへ移動する（\Cref{tab:layout}）．

\begin{table}[htbp]
  \centering
  \caption{実験空間の主要座標}
  \label{tab:layout}
  \begin{tabular}{ll}
    \toprule
    名称 & 座標 $(x, y, z)$ \\
    \midrule
    Station A（開始） & $(22,\, 0.2,\, 9)$ \\
    Station B（ゴール） & $(9,\, 0.2,\, 58)$ \\
    AGV床面高 & $y = \SI{0.15}{m}$ \\
    \bottomrule
  \end{tabular}
\end{table}

\subsection{被験者とデザイン}

\begin{itemize}[leftmargin=2em]
  \item \textbf{デザイン}: 被験者内計画，3条件（Baseline / No-AR / Proposed）
  \item \textbf{反復}: 各条件10ケース（計30試行／被験者）
  \item \textbf{ケース固定}: 同一ケース番号ではAGV台数（10--20），速度（\SIrange{1.0}{2.0}{m/s}），
        軌道をシード固定し，条件間でAGV挙動を一致させる
  \item \textbf{被験者数}: 予備実験後に決定（パイロット $n \ge 8$ を推奨）
\end{itemize}

\subsection{手続き}

\begin{enumerate}[leftmargin=2em]
  \item 説明と同意取得，HMD装着，操作練習
  \item 実験者がPC上のメニューから条件とケース（1--10）を選択
  \item 3秒のイントロ表示後，Station Aにテレポート
  \item 被験者はStation B（半径\SI{3}{m}以内）へ移動
  \item ゴール到着で試行終了，データ保存，Station Aへ復帰
  \item 全条件・全ケース完了まで反復
\end{enumerate}

\subsection{操作}

\begin{itemize}[leftmargin=2em]
  \item \textbf{移動}: 左スティック（HMDの向き基準）
  \item \textbf{視点}: HMDの頭部追従
  \item \textbf{歩行速度}: \SI{1.4}{m/s}
\end{itemize}

\subsection{従属変数}

\begin{table}[htbp]
  \centering
  \caption{従属変数と操作定義}
  \label{tab:dv}
  \begin{tabular}{llp{5.5cm}}
    \toprule
    変数 & 単位 & 定義 \\
    \midrule
    完了時間 & s & Station A出発からStation B到達まで \\
    ニアミス回数 & 回 & AGV半径\SI{0.5}{m} + 被験者半径\SI{0.3}{m}
    = \SI{0.8}{m} 以内接近（AGVごとクールダウン\SI{1.5}{s}） \\
    移動距離 & m & 水平面上の軌跡長（\SI{0.5}{s}間隔サンプリング） \\
    \bottomrule
  \end{tabular}
\end{table}

\subsection{データ記録}

1 PlayセッションにつきExcel（.xlsx）1ファイルを出力する．
Summaryシートに条件，ケース番号，従属変数を記録し，
ケース別シートに時系列位置 $(x,y,z)$ を保存する．

\section{実験の妥当性}
\label{sec:validity}

本節では，計画したVR実験が研究課題に答えるに足る\textbf{妥当性}を，
Campbell \& Stanley の枠組みと現代HRI研究の慣行に沿って論じる．

\subsection{構成概念の妥当性（Construct Validity）}

\textbf{測定したい概念}と\textbf{操作・指標}の対応関係を次のように定義する．

\begin{table}[htbp]
  \centering
  \caption{構成概念と操作化}
  \label{tab:construct}
  \begin{tabular}{ll}
    \toprule
    構成概念 & 操作化 \\
    \midrule
    移動効率 & 完了時間，移動距離 \\
    安全性（ニアミス） & 閾値\SI{0.8}{m}以内の接近回数 \\
    経路情報の提示 & Baseline / No-AR / Proposed \\
    危険度連動表示 & Proposed条件のみ $R$ に基づく色・不透明度 \\
    AGV環境の複雑さ & 10--20台，ピックアップ／ドロップ反復 \\
    \bottomrule
  \end{tabular}
\end{table}

Proposed条件では，TTCと近接距離という\textbf{客観的指標}から $R$ を算出し，
色相・不透明度へ決定論的に写像するため，
操作者バイアスが視覚刺激に混入しにくい．
また，Baseline条件は「情報量一定・危険度エンコードなし」，
No-ARは「情報ゼロ」として，\textbf{危険度連動の効果}を分離して評価できる．

\subsection{内部妥当性（Internal Validity）}

条件間比較の信頼性を高めるため，以下の統制を行う．

\begin{itemize}[leftmargin=2em]
  \item \textbf{AGV挙動の固定}:
        同一ケース番号では乱数シードを固定し，
        3条件でAGV台数・速度・経路計画を一致させる
  \item \textbf{空間・タスクの固定}:
        Station A/B，NavMesh，工場レイアウトは全条件共通
  \item \textbf{被験者内デザイン}:
        個人差（歩行速度，空間認知能力）を条件効果から分離
  \item \textbf{自動計測}:
        完了時間，ニアミス，軌跡はプログラムで記録し，観察者バイアスを低減
\end{itemize}

\textbf{残存する脅威}として，(1)条件提示順序の学習効果，
(2)VR酔い，(3)被験者の危険感の個人差が挙げられる．
これらに対し，条件順序のカウンターバランス（またはラテン方格），
休憩の挿入，事前練習を計画する．

\subsection{外部妥当性（External Validity）}

VR工場環境は実工場の完全な再現ではないが，
\textbf{複数AGVと作業員の通路共有}という本質的状況を再現する．
AGVはNavMesh上を走行し，被験者はCharacterControllerで歩行するため，
\textbf{衝突回避の時間的圧力}は現実に近い形で生じる．

一方，視覚・聴覚（エンジン音等），物理的触覚は簡略化されている．
したがって本実験の結果は，
「\textbf{AR経路eHMIの相対比較}」として解釈するのが適切であり，
絶対的な事故率予測には追加の現場検証が必要である．

\subsection{統計的結論の妥当性}

\begin{itemize}[leftmargin=2em]
  \item 従属変数ごとに正規性・球面性を確認
  \item 主分析: 3条件の反復測定ANOVA（または非パラメトリック代替）
  \item 事後検定: Proposed vs.\ Baseline，Proposed vs.\ No-AR（Bonferroni補正）
  \item 効果量: $\eta^2$ または Cohen's $d$ を報告
  \item 有意水準: $\alpha = 0.05$
\end{itemize}

\subsection{実験計画の妥当性まとめ}

\begin{table}[htbp]
  \centering
  \caption{妥当性の総合評価}
  \label{tab:validity_summary}
  \begin{tabular}{p{3.2cm}p{9.5cm}}
    \toprule
    観点 & 評価 \\
    \midrule
    研究課題との整合 & RQ（危険度連動表示の効果）に直接対応した3条件比較 \\
    操作化 & $R$ の数理定義と実装が一致（再現可能） \\
    統制 & シード固定によるAGV挙動一致 \\
    一般化 & VR内比較として妥当；現場一般化は今後の課題 \\
    \bottomrule
  \end{tabular}
\end{table}

本実験は，\textbf{「危険度連動eHMIが，情報なし・均一表示と比べて
効率と安全のトレードオフを改善するか」}を検証するうえで，
方法論的に妥当な計画であると言える．

\section{実装概要}
\label{sec:implementation}

提案手法はUnity 6（URP）プロジェクト \texttt{HRI-Robot} に実装済みである．
\Cref{tab:impl}に論文概念と実装クラスの対応を示す．

\begin{table}[htbp]
  \centering
  \caption{論文概念と実装の対応}
  \label{tab:impl}
  \begin{tabular}{ll}
    \toprule
    論文概念 & 実装クラス \\
    \midrule
    動的基準線 & \texttt{DynamicCrossingLineTracker} \\
    危険度スコア & \texttt{VehicleRiskCalculator} \\
    色・不透明度 & \texttt{RiskToVisualMapper} \\
    経路描画 & \texttt{PathRenderer} \\
    条件切替 & \texttt{AGVPathVisualizer} \\
    実験制御 & \texttt{ExperimentManager} \\
    データ出力 & \texttt{MeasurementHub}, \texttt{ExperimentWorkbook} \\
    \bottomrule
  \end{tabular}
\end{table}

\subsection{AGVシミュレーション}

各AGVは \texttt{MovingToPickup} $\rightarrow$ \texttt{DwellAtPickup}
$\rightarrow$ \texttt{MovingToDrop} $\rightarrow$ \texttt{DwellAtDrop}
のサイクルを反復する．
経路計画はNavMesh上のA*相当処理（\texttt{AGVRoutePlanner}）で，
床面高 $y=\SI{0.15}{m}$ に固定する．
全AGVは\SI{0.1}{s}固定クロックで一括更新される（\texttt{AGVFleetOrchestrator}）．

\subsection{実験フロー}

\texttt{ExperimentStartMenu} により条件（Baseline / No-AR / Proposed）と
ケース1--10を選択する．
ケース開始時に \texttt{AGVSpawner.RestartWithSeed(seed)} を呼び，
ゴール到着時にExcelへ結果を追記する．

\subsection{再現性}

同一リポジトリ，同一シーン，同一シード列により，
他研究者が実験刺激を再現可能である．
倉庫3Dアセット（Unity Warehouse）はローカルインポートとし，
GitHubには含めない（ライセンス・容量のため）．

\section{考察（予定）}
\label{sec:discussion}

\subsection{期待される結果パターン}

\begin{itemize}[leftmargin=2em]
  \item \textbf{Proposed $<$ No-AR（安全）}:
        危険AGVの経路が赤く強調され，回避が促進される
  \item \textbf{Proposed $<$ Baseline（効率）}:
        非危険AGVの表示が抑制され，注意資源が節約される
  \item \textbf{Proposed $>$ Baseline（安全）}:
        均一表示では見落としていた危険AGVが強調される
\end{itemize}

\subsection{設計上の論点}

\begin{enumerate}[leftmargin=2em]
  \item \textbf{視線追従基準線}:
        目的地を仮定しないため，作業中の視線変化に頑健である一方，
        視線がAGV方向を向いていないときの検出限界がある．
  \item \textbf{2D交差判定}:
        床面走行AGVを前提としたXZ平面近似であり，
        段差・立体障害は今後の拡張課題である．
  \item \textbf{不透明度の量子化}:
        連続色と段階的不透明度の組合せが，
        被験者の危険認知に与える影響は主観評価で補完可能である．
\end{enumerate}

\subsection{限界と今後の展開}

\begin{itemize}[leftmargin=2em]
  \item 現場環境（騒音，荷物，疲労）の再現
  \item 被験者数の拡大と現場パイロット
  \item 基準線半幅 $w$，$T_\text{max}$ の感度分析
  \item 主観的ワークロード（NASA-TLX）・信頼尺度の追加
\end{itemize}

\section{おわりに}
\label{sec:conclusion}

本研究は，工場内AGV環境において，
残存走行経路をAR風に可視化するeHMIに\textbf{危険度連動の視覚エンコーディング}を
導入する提案手法を設計した．
危険度は，視線方向の動的基準線とAGV残存経路のTTC，
および近接距離から統合的に算出される．

VR実験（被験者内3条件×10ケース）により，
移動効率とニアミスを定量的に比較し，
均一表示（Baseline）および非表示（No-AR）に対する
提案手法の優位性を検証する計画である．
シード固定によるAGV挙動統制，自動計測，構成概念の明示的操作化により，
実験計画は方法論的に妥当である．

今後，データ収集と統計分析を行い，
スマートファクトリーにおける人--AGV共存の安全HRI設計に
知見を還元する．


\bibliographystyle{jplain}
\begin{thebibliography}{99}
  \bibitem{matthiesen2021survey}
  S.~Matthiesen et al.,
  ``Survey on eHMI concepts for autonomous vehicles,''
  \textit{IEEE Transactions on Intelligent Transportation Systems}, 2021.
  % TODO: 正式な文献情報に差し替え

  \bibitem{francis2020eHMI}
  C.~Francis et al.,
  ``A review of external human-machine interfaces for autonomous vehicles,''
  \textit{Proceedings of the Human Factors and Ergonomics Society Annual Meeting}, 2020.
  % TODO: 正式な文献情報に差し替え

  \bibitem{markkula2020}
  A.~Markkula et al.,
  ``Explaining human interactions with external vehicle displays,''
  \textit{Transportation Research Part F}, 2020.
  % TODO: 正式な文献情報に差し替え
\end{thebibliography}

\end{document}
